\section{Conclusion}
\label{sec:conclusion}

End-to-end encryption protects an attachment while it crosses an untrusted
relay, but it cannot determine which decrypted representation an endpoint
should expose to a parser or renderer. \AFD\ addresses this boundary through a
composed endpoint-local method. An outgoing object is processed before
encryption and an incoming object after decryption. The relay remains blind.
Exactly three routing sources select a constrained handler, namely the
normalized portable name, caller-supplied media type, and bounded byte evidence.
An independent complete output parser participates in delivery authorization.
Policy-indexed reconstruction strength prevents weaker processing from
satisfying a stronger requirement. Publication is clean-only and independently
validated. A pre-specified evidence contract fixes the empirical obligations
before values are admitted. These five elements, composed as a user-benefiting
endpoint defense in an encrypted workflow, define the residual novelty.

The method's scientific value rests on the distinctions it preserves. A
structural reconstruction claim is not presented as semantic disarm. A
successful decoder is not presented as complete format validation. A carrier
scenario associated with a virus, worm, Trojan, ransomware, or spyware label is
not presented as malware classification. A downloader, dropper, or stager role
is not inferred from an unsupported carrier. Passive metadata is treated as
privacy leakage, not spyware. Fuzzing is not presented as proof. Planned server
measurements are not presented as completed results.

The implementation is organized so that the construction can be traced through
versioned policy, media handlers, independent validators, bounded child
processes, verdict derivation, and publication interfaces. The declared format
profiles cover static raster images, GIF, APNG, animated WebP, MP3, WAV, FLAC, Ogg,
constrained ISO BMFF, and constrained WebM with explicitly different
reconstruction strengths. Unsupported active classes are assigned a blocking
decision by the evaluated media-only profile.

The extended evaluation protocol makes the remaining empirical obligation
precise. The pre-specified evidence contract freezes the software, policy,
corpus, environment, mutations, properties, fault campaign, statistical
procedures, and release gates before numbers are admitted. Once that protocol
is executed on the controlled server,
the resulting evidence can establish profile-specific effectiveness, usability,
fault behavior, and cost for the pinned release. Until then, the contribution
is a claim-bounded methodology and an inspectable implementation, not a
completed population-level assertion of malware prevention.
